% from ega4-18-fr.tex
\oldpage[IV-4]{134}

qui est aussi la topologie $\mathfrak s$-préadique, où $\mathfrak s$ est le radical de l'anneau semi-local $B$, car $B/\mathfrak rB$ est un anneau artinien de radical $\mathfrak s/\mathfrak rB$. On sait alors (Bourbaki, \emph{Alg. comm.}, chap. III, \textsection~2, n\textsuperscript{o}~13, cor. de la prop. 19) que $B$ est composé direct d'anneaux locaux.

\begin{proposition}[18.5.15]
\label{IV.18.5.15}
Soient $A$ un anneau hensélien, $\mathfrak r$ son radical. Alors le foncteur $B\longmapsto B/\mathfrak rB$ est une équivalence de la catégorie des $A$-algèbres finies et étales avec la catégorie des $(A/\mathfrak r)$-algèbres finies et étales.
\end{proposition}

Le fait que le foncteur de l'énoncé soit pleinement fidèle est un cas particulier de (18.5.12). Pour montrer que ce foncteur est une équivalence de catégories, on peut se borner au cas où $A$ est local; il suffit alors d'appliquer (18.1.1) (pour les morphismes étales), à $S=\Spec(A)$ et $S_0=\Spec(A/\mathfrak r)$, réduit à un seul point.

\begin{remarks}[18.5.16]
\label{IV.18.5.16}
\begin{enumerate}
  \item[(i)] Nous ignorons si la proposition (18.5.15) se généralise à un couple hensélien $(S,S_0)$, même lorsque $S$ est affine et noethérien.
  \item[(ii)] Soient $A$ un anneau noethérien, $\mathfrak J$ un idéal de $A$ tel que $A$ soit \emph{séparé et complet} pour la topologie $\mathfrak J$-préadique. Alors le couple $(\Spec(A),\Spec(A/\mathfrak J))$ est hensélien: en effet, pour toute $A$-algèbre finie $B$, $B$ est séparé et complet pour la topologie $\mathfrak J$-préadique ($0_{\mathrm I}$, 7.3.6). Remplaçant $B$ par $A$ et $\mathfrak J B$ par $\mathfrak J$, tout revient donc à voir que l'application qui, à tout idempotent de $A$, fait correspondre sa classe mod. $\mathfrak J$, est \emph{bijective}. Or on a $A=\varprojlim (A/\mathfrak J^n)$. Notons $\operatorname{Idem}(A)$ l'ensemble des idempotents de $A$, et pour tout homomorphisme d'anneaux $\varphi:A\to B$, soit $\operatorname{Idem}(\varphi)$ l'application de $\operatorname{Idem}(A)$ dans $\operatorname{Idem}(B)$ restriction de $\varphi$; il résulte de la définition de la limite projective que l'on a
  \[
    \operatorname{Idem}(A)=\varprojlim_n\operatorname{Idem}(A/\mathfrak J^n)
  \]
  pour les applications $\psi_{nm}:\operatorname{Idem}(A/\mathfrak J^m)\to\operatorname{Idem}(A/\mathfrak J^n)$ restrictions des applications canoniques $A/\mathfrak J^m\to A/\mathfrak J^n$. Mais puisque $\Spec(A/\mathfrak J^n)\to\Spec(A/\mathfrak J^m)$ est un homéomorphisme, les $\psi_{nm}$ sont des bijections (comme on l'a vu dans la démonstration de (18.5.3)); cela prouve donc notre assertion.
\end{enumerate}
\end{remarks}

\begin{theorem}[18.5.17]
\label{IV.18.5.17}
Soient $A$ un anneau local hensélien, $S=\Spec(A)$, $s$ le point fermé de $S$. Pour tout morphisme lisse $f:X\to S$, l'application canonique
\[
  \Gamma(X/S)\longrightarrow\Gamma(X_s/k(s))
\]
(où $X_s=f^{-1}(s)$) est surjective.
\end{theorem}

La donnée d'une $k(s)$-section de $X_s$ équivaut à celle d'un point $x\in X$ au-dessus de $s$ rationnel sur $k(s)$, et il s'agit de prouver qu'il existe une $S$-section $u:S\to X$ telle que $u(s)=x$. Compte tenu de (17.16.3, (i)), on peut supposer que $f$ est étale. Alors la conclusion résulte du critère (18.5.11, $b)$). (Le lecteur notera qu'en vertu de ce critère, la validité de (18.5.17) pour un anneau local donné $A$ est nécessaire et suffisante pour que $A$ soit hensélien).

\begin{remark}[18.5.18]
\label{IV.18.5.18}
Les conditions de (18.5.11) sont encore équivalentes à la suivante:
\begin{enumerate}
  \item[d)] Pour tout morphisme localement de type fini $f:X\to S$ et tout point $x\in X$ tel que $f(x)$ soit le point fermé $s$ de $S$ et que $\sh{O}_{X,x}/\mathfrak m_s\sh{O}_{X,x}$ soit un corps canoniquement isomorphe à $k(s)=k$, il existe un voisinage ouvert $U$ de $x$ dans $X$ tel que $f|U$ soit une immersion fermée.
\end{enumerate}
\end{remark}

